\chapter{Building \linnet{}}
\label{secBuildingLinNet}

\linnet{} is distributed as source code together with makefiles to compile
and link the code. The distribution contains ready-to-use binaries for
Windows, as debug and production compilation\footnote{Windows 11 has been
used.}. In all other environments a build from the sources needs to be
done prior to the use of \linnet{}.

To start the build you need to open a shell window and \code{cd} into the
root directory of the cloned project; this is the folder, where the file
\file{GNUmakefile} is located. Here, you issue the command:
\begin{verbatim}
make -sO build
\end{verbatim}
The makefile has some more useful targets and options than just
\ident{build}. Type \code{make help} to find out.


\section{Portability of makefiles}

The makefile -- actually it is a set of nested makefiles -- is compatible
with the GNU make processor of at least revision 3.81; 3.80 is not
sufficient.\footnote{Latest tests have been run with GNU make 4.4.1 and it
is recommended to use a recent revision of GNU make.} On a Windows
platform, it needs to be the MinGW port of GNU make; the Cygwin port has
another shell interface and doesn't work.

Different derivates of make always tended to be quite incompatible to one
another and other derivates of make will hence require heavy modifications
of the makefiles.

The makefile directly supports the build of the software for Windows with
GCC in the MinGW 32 and 64 Bit ports. The compilation and the compiled
software have been tested with:

\begin{itemize}
  \item mingw-i686-8.1.0-release-win32-dwarf-rt\_v6-rev0 (32 Bit, Windows 11)
  \item mingw-w64-x86\_64-8.1.0-posix-seh-rt\_v6-rev0 (64 Bit, Windows 11)
  \item mingw-x86\_64-11.2.0-release-win32-sjlj-rt\_v9-rev1 (64 Bit, Windows 11)
  \item GCC 12.2.0, x86\_64-posix-seh-rev2 (64 Bit, Windows 11)
  \item mingw-w64-gcc-13.2.0-x64 (64 Bit, Windows 11)
  \item mingw-w64-gcc-mcf\_20260501\_16.1.0\_x86-msvcrt (32 Bit, Windows 11)
  \item mingw-w64-gcc-mcf\_20260501\_16.1.0\_x64-msvcrt (64 Bit, Windows 11)
\end{itemize}

\noindent
The build depends on two packages:

\begin{itemize}
  \item The GNU compiler collection GCC (see above).
  \item A set of common UNIX command line tools like \code{echo},
    \code{mkdir} and \code{rm}. On a Windows platform, these tools aren't
    natively available. You will probably find them inside the
    installation directory of your Git 
    (\file{\textless{}gitinstalldir\textgreater{}\textbackslash{}usr\textbackslash{}bin})
    or you can download them from the MinGW project.
\end{itemize}
 
The makefile needs to know, where these packages are installed. Two
options are supported. 

You can add the directories with the binaries of GCC and the UNIX tools
to your global search path, i.e., ensure that the environment variable
\ident{PATH} contains these directories.

Or you set the following two environment variables, which are evaluated by
the makefile:\footnote{Although designed for Windows, where typically
several different competing ports of GCC can reside, the environment
variables can be used for Linux and Mac OS, too.}

\begin{itemize}
  \item \ident{CC\_BIN}: This variable explicitly specifies the folder
    with the compiler binaries, i.e., the folder, where \file{gcc.exe} is
    located in.
  \item \ident{UNIX\_TOOLS\_BIN}: This variable explicitly specifies the
    directory with the UNIX tools, i.e.,
    \file{\%UNIX\_TOOLS\_BIN\%\textbackslash{}echo.exe} or
    \file{\$(UNIX\_TOOLS\_BIN)/echo} should be a valid executable file.
\end{itemize}

An advantage of the second method is that the variables can be set
still on the command line of the make processor. This allows, e.g., to
switch on the fly between 32 Bit and 64 Bit builds, and it overrules a
persistent setting made in the global system search path variable
\ident{PATH}.

The build process has been designed to support the compilation under Linux
and Mac OS, too. GNU make's \code{if}/\code{else}/\code{endif} statements
have been applied to do so. Unfortunately, this has been barely tested and
might not work out of the box. In case of doubts, you should have a look
at the tool localization in \file{locateTools.mk} and at the compiler and
linker command line options in \file{compileAndLink.mk}.
\file{locateTools.mk} offers some debugging support; look for commented
\code{\$(info ...)} statements and uncomment them temporarily.

\subsection{Resource compiler}

Under Windows the application icons are added to the executable binary
file. This has no functional aspect; the image data is not accessed from
the functional code. It's just a gimmick that allows to create illustrated
associations between \linnet{} in- and output files and the executable in
the file system browser.

In the makefile the according build rules are placed in
\code{if}/\code{endif} clauses; on a Linux or Mac OS system, they have no
effect. Similar mechanisms will probably exist for these operating systems
as well. You may consider to extend the makefiles by adding appropriate
\code{else} paths; an OS indicating variable already exists.


\section{Portability of source code}

The source code itself is platform independent; only the abstract
functions from the GNU C library are used. The language standard is C99.

\subsection{GNU extensions}

In module \ident{rat\_rationalNumber} a GNU extension is used: Initialized
structs are used as RHS of assignment expressions. This has been done for
convenience only, to achieve a tense and meaningful representation of the
required functionality. It is easy and straightforward to eliminate these
(few) constructs. All the rest is according to the language standard C99.

\subsection{Incompatibilities of linked libraries}

Far the most necessary code modifications will result from
incompatibilities of the linked libraries. Most obvious is the use of
GCC's Basic Program/System Interface, the command line evaluation support.
All of the command line evaluation related code will probably require a
complete re-implementation with another compiler tool chain. Please refer
to module \ident{opt\_getOpt} for details.

The standard function snprintf was found to be buggy in elder releases of
GCC and up to 8.1. A workaround has been implemented in file
\file{snprintf.h}, which is compiled only for explicitly selected compiler
releases. In debug compilation, an assertion checks if the problem is
present or not. Code execution is halted if your compiler would require
the workaround. A related source code comment tells, what to do in this
situation.

\section{Re-entrance of code}

The implementation of \linnet{} is widely but not completely reentrant.
This can become an issue if the code should be integrated into a
concurrent environment; for example a client-server structure is set up or
multi-core support is implemented to speed up the computations. As long as
the code is used as of today in a single-threaded self-contained
application the following considerations are irrelevant.

Although the implementation of module \ident{log\_logger} is made in a
reentrant style, the logging process is not reentrant by principle. All
logger instances write to the same global stream \ident{stdout} and
concurrent use would lead to a mess of text at the console. Furthermore,
all modules of \linnet{} use one and the same global logger instance and
they hence share the same file stream, too; here, we'd expect a similar
mess.

For module \ident{pci\_parserCircuit}, the parser of the circuit netlist
files, the decision to use a few global variables (e.g., for error
reporting by side effect) has been taken just for convenience, to keep the
code lean. An application for a fully reentrant implementation was simply
not in the scope of the development. It's however straightforward and only
little effort to make the parser reentrant. One just has to put the few
global variables into a struct, which gets the meaning of a parser
instance, add a constructor/destructor pair and pass the pointer to the
parser instance to all the functions.

The situation differs for the other affected module
\ident{rat\_rationalNumber}, the implementation of an exact arithmetics
for rational numbers. Here, the decision to use global data and to loose
re-entrance has been taken intentionally for performance reasons. A global
variable is used to report an overflow. In addition, a global logger
instance is used to visibly report problems. This logger isn't reentrant,
neither. The client code can check the global variable after completing
the operations. The flag is persistent and can thus be checked after a
bunch of operations as an overall result. To avoid the flag we could
return the overflow indication with every operation but this would lead to
a manifold of computational effort. Or we could introduce rational number
processor instances having each their own individual error handlers. The
overhead of passing the instance pointer to all the functions would
probably be less than in the first suggestion but still significant.

If the software is used in an environment, where re-entrance is a must then
serialized access to the module could also be an option. A client is
queued until the other one has completed its (bunch of) operations and has
checked the global error information.